\subsection{QC-007: exact centered determinant samples for the invariant comparison}

\paragraph{Statement QC-007.}
The file \texttt{centered\_dense\_830.json} gives $830$ explicitly
reproducible points of the centered nine-by-nine determinantal
parameter image modulo $101$, with the $161$ normalized
symmetric-tensor coordinates of orders $2$ through $8$ at each
point. Every source pencil has an integer lift with constant
coefficient $I_9$ and exactly traceless remaining matrices.
This is evaluation data for an isotypic rank comparison, and
does not by itself assert a rank or a multiplicity obstruction.

\begin{proof}
Choose three integer matrices $A_1,A_2,A_3$ of size nine and set
\[
 F(x)=\det(x_0I_9+x_1A_1+x_2A_2+x_3A_3).
 \tag{QC-007.1}
\]
The generator starts with $s=20261122$ and uses
$s\leftarrow1664525s+1013904223\pmod {2^{32}}$ and
$c=(s\mathbin{\mathrm{div}}2^{16})\bmod101$. It fills each
matrix in row order, including all $81$ draws, and then
replaces its last diagonal entry by minus the sum of its
first eight diagonal entries. There is no reduction modulo
$101$ in this last assignment, so each displayed matrix has
trace zero over $\mathbb Z$. Three successive matrices form
one sample; the state is retained for the next sample.
Consequently $[x_0^9]F=1$ and
$[x_0^8x_i]F=\operatorname{tr}(A_i)=0$ exactly. Thus the
ordinary centering change is already the identity at all
sample points.

For clarity, write
\[
 F=x_0^9+\sum_{j=2}^9\binom9j x_0^{9-j}
 T_j(x_1,x_2,x_3),
 \qquad
 T_j(x)=\sum_{a+b+c=j}\frac{j!}{a!b!c!}
 T_{abc}\,x_1^ax_2^bx_3^c.
 \tag{QC-007.2}
\]
The normalized tensor entries are therefore
\[
 T_{abc}=
 [x_0^{9-j}x_1^ax_2^bx_3^c]F\,
 \frac{(9-j)!a!b!c!}{9!}.
 \tag{QC-007.3}
\]
All denominators are invertible modulo $101$. The stored
coordinates are precisely (QC-007.3), in the explicit order
$(j,a,b,j-a-b)$ with $j=2,\ldots,8$, $a=0,\ldots,j$ and
$b=0,\ldots,j-a$. There are
$\sum_{j=2}^{8}\binom{j+2}{2}=161$ entries.

Here is an exact coefficient-reconstruction algorithm requiring
only scalar Gaussian determinants. Let
$f(v)=F(1,v_1,v_2,v_3)$ and evaluate it at all $220$ nodes
$\alpha\in\mathbb Z_{\geq0}^3$ with $|\alpha|\leq9$.
For every such $\alpha$, compute
\[
 d_\alpha=
 \sum_{\beta\leq\alpha}(-1)^{|\alpha-\beta|}
 \prod_{i=1}^3\binom{\alpha_i}{\beta_i}\,f(\beta).
 \tag{QC-007.4}
\]
The multivariate Newton identity is
\[
 f(v)=\sum_{|\alpha|\leq9}
 d_\alpha\prod_{i=1}^3\binom{v_i}{\alpha_i}.
 \tag{QC-007.5}
\]
Indeed the binomial polynomials form a basis of polynomials
of total degree at most nine, and the forward-difference
operators in the three coordinates extract their coefficients.
Equivalently, their values on these nodes form a triangular
matrix with diagonal one in any ordering by total degree.
Their factorial denominators are invertible because $101>9$.
Expanding each falling factorial in (QC-007.5) recovers every
coefficient of $f$ exactly. No probabilistic interpolation
assumption is involved.

The standard-library program \texttt{centered\_det9\_sampler.py}
implements (QC-007.4)--(QC-007.5), and
\texttt{generate\_centered\_dataset.py} writes the data set.
Each sample also has three explicit points at which a direct
Gaussian determinant equals the reconstructed polynomial;
all $2490$ comparisons passed. These additional evaluations
are consistency checks, while (QC-007.4)--(QC-007.5) are the
coefficient certificate itself. As a separate implementation
check, \texttt{check\_centered\_sampler.py} compared all
coefficients against direct six-term Leibniz expansion on
$24$ three-by-three pencils and against multiplication of
three such expansions on $12$ block diagonal nine-by-nine
pencils. All $3120$ coefficient comparisons passed.

The tensor values, one byte per residue in sample-major and
then coordinate-major order, have SHA-256
\begin{quote}\ttfamily\small
acff669d4dc6b8961c6d267567536b18fa7e4252b2e9805365efba381632e6f4.
\end{quote}
The complete data file has SHA-256
\begin{quote}\ttfamily\small
6b61b4279eabf95811b22d1d366c0982a265264e9205d25226eb570017cad536.
\end{quote}
The adjacent manifest specifies both byte encodings and the
source-matrix digest. All sample states, sanity points, and
the first source pencil are also stored. Re-running the
generator reproduces the entire file; the runtime is confined
to the separate manifest and does not affect its digest.
\end{proof}

\paragraph{Independent review of the invariant-basis construction.}
The construction in \texttt{bracket\_basis.py} fixes the first
of eight three-dimensional epsilon tensors to its identity
orientation. For an orbit $O$ of tensor monomials under a
global permutation of the three coordinate colors, let $q_O$
be the sum of signed reduced-orientation terms whose monomial
lies anywhere in $O$. Restoring all six orientations gives
coefficient
\[
 \frac{6}{|O|}\,q_O
 \tag{QC-007.6}
\]
at each member of $O$. To see this, a global permutation
$\sigma$ multiplies the product of epsilon signs by
$\operatorname{sgn}(\sigma)^8=1$, and each target member of
$O$ is reached $|\operatorname{Stab}(O)|=6/|O|$ times.
This accounts for the complete integral normalization.
No further multinomial factors occur when tensor entries
are the normalized entries (QC-007.3).

The code's rank calculation uses $q_O$ instead of the full
coefficients. This only rescales each coefficient column by
$6/|O|\in\{1,2,3,6\}$, a unit modulo $101$, so it preserves
rank. The separate standard-library verifier restores all
full coefficients and checks the integral raising identity.
The map on a tensor coordinate indexed by $(a,b,c)$ has
coefficient $b$ and target $(a+1,b-1,c)$; the verifier
also includes the multiplicity of that variable in the
commutative tensor monomial. Global color symmetry conjugates
this identity to the other simple raising identity. Together
with weight $(8,8,8)$, these are the asserted characteristic-zero
invariance conditions. The whole expansion, index mapping,
orbit restoration, and verifier code were read independently;
no fault was found. This review does not replace the complete
numerical basis verification owned by the support branch.
